\documentclass[12pt]{article}

% Idioma y codificación
\usepackage[spanish, es-tabla, es-notilde]{babel}       %es-tabla para que se titule "Tabla"
\usepackage[utf8]{inputenc}

% Márgenes
\usepackage[a4paper,top=3cm,bottom=2.5cm,left=3cm,right=3cm]{geometry}

% Comentarios de bloque
\usepackage{verbatim}

% Paquetes de links
\usepackage[hidelinks]{hyperref}    % Permite enlaces
\usepackage{url}                    % redirecciona a la web

% Más opciones para enumeraciones
\usepackage{enumitem}

% Personalizar la portada
\usepackage{titling}

% Paquetes de tablas
\usepackage{multirow}

% Para añadir el símbolo de euro
\usepackage{eurosym}

% Paquete para subrayar con diferentes estilos
\usepackage[normalem]{ulem}


%------------------------------------------------------------------------

%Paquetes de figuras
\usepackage{caption}
\usepackage{subcaption} % Figuras al lado de otras
\usepackage{float}      % Poner figuras en el sitio indicado H.


% Paquetes de imágenes
\usepackage{graphicx}       % Paquete para añadir imágenes
\usepackage{transparent}    % Para manejar la opacidad de las figuras

% Paquete para usar colores
\usepackage[dvipsnames, table, xcdraw]{xcolor}
\usepackage{pagecolor}      % Para cambiar el color de la página

% Habilita tamaños de fuente mayores
\usepackage{fix-cm}

% Para los gráficos
\usepackage{tikz}
\usepackage{forest}

% Para poder situar los nodos en los grafos
\usetikzlibrary{positioning}

% Para evitar problemas entre babel y tikz
\usetikzlibrary{babel}


%------------------------------------------------------------------------

% Paquetes de matemáticas
\usepackage{mathtools, amsfonts, amssymb, mathrsfs}
\usepackage[makeroom]{cancel}     % Simplificar tachando
\usepackage{polynom}    % Divisiones y Ruffini
\usepackage{units} % Para poner fracciones diagonales con \nicefrac

\usepackage{pgfplots}   %Representar funciones
\pgfplotsset{compat=1.18}  % Versión 1.18

\usepackage{tikz-cd}    % Para usar diagramas de composiciones
\usetikzlibrary{calc}   % Para usar cálculo de coordenadas en tikz

%Definición de teoremas, etc.
\usepackage{amsthm}
%\swapnumbers   % Intercambia la posición del texto y de la numeración

\theoremstyle{plain}

\makeatletter
\@ifclassloaded{article}{
  % --- Configuración para ARTICLE ---
  \newtheorem{teo}{Teorema}[section]
  \theoremstyle{definition}
  \newtheorem{definicion}{Definición}[section]
  \newtheorem{ejercicio}{Ejercicio}[section]

  \renewcommand{\theteo}{%
    \ifnum\value{section}=0 \arabic{teo}\else \thesection.\arabic{teo}\fi
  }
  \renewcommand{\thedefinicion}{%
    \ifnum\value{section}=0 \arabic{definicion}\else \thesection.\arabic{definicion}\fi
  }
  \renewcommand{\theejercicio}{%
    \ifnum\value{section}=0 \arabic{ejercicio}\else \thesection.\arabic{ejercicio}\fi
  }
}{
  % --- Configuración para BOOK / REPORT ---
  \newtheorem{teo}{Teorema}[chapter]
  \theoremstyle{definition}
  \newtheorem{definicion}{Definición}[chapter]
  \newtheorem{ejercicio}{Ejercicio}[section]

  % Si no hay sección: Capítulo.Teorema (ej. 2.1)
  % Si hay sección:    Capítulo.Sección.Teorema (ej. 2.3.1)
  \renewcommand{\theteo}{%
    \ifnum\value{section}=0
      \thechapter.\arabic{teo}%
    \else
      \thesection.\arabic{teo}% (\thesection ya contiene el número de capítulo)
    \fi
  }
  \renewcommand{\thedefinicion}{%
    \ifnum\value{section}=0
      \thechapter.\arabic{definicion}%
    \else
      \thesection.\arabic{definicion}%
    \fi
  }
  \renewcommand{\theejercicio}{%
    \ifnum\value{section}=0
      \arabic{ejercicio}%
    \else
      \thesection.\arabic{ejercicio}% Formato por defecto: ej. 2.2.1
    \fi
  }
}
\makeatother

% Teoremas dependientes o con contador compartido
\newtheorem{coro}{Corolario}[teo]
\newtheorem{prop}[teo]{Proposición}
\newtheorem{lema}[teo]{Lema}
\newtheorem*{lema*}{Lema}

\theoremstyle{remark}
\newtheorem*{observacion}{Observación}

\theoremstyle{definition}
\newtheorem*{notacion}{Notación}
\newtheorem*{ejemplo}{Ejemplo}
\newtheorem*{ejercicio*}{Ejercicio}


% \renewcommand\qedsymbol{$\blacksquare$}         % Cambiar símbolo QED
%------------------------------------------------------------------------

% Paquetes para encabezados
\usepackage{fancyhdr}
\pagestyle{fancy}
\fancyhf{}

\newcommand{\helv}{ % Modificación tamaño de letra
\fontfamily{}\fontsize{12}{12}\selectfont}
\setlength{\headheight}{15pt} % Amplía el tamaño del índice


%\usepackage{lastpage}   % Referenciar última pag   \pageref{LastPage}
%\fancyfoot[C]{%
%  \begin{minipage}{\textwidth}
%    \centering
%    ~\\
%    \thepage\\
%    \href{https://losdeldgiim.github.io/}{\texttt{\footnotesize losdeldgiim.github.io}}
%  \end{minipage}
%}
\fancyfoot[C]{\thepage}
\fancyfoot[R]{\href{https://losdeldgiim.github.io/}{\texttt{\footnotesize losdeldgiim.github.io}}}

%------------------------------------------------------------------------

% Conseguir que no ponga "Capítulo 1". Sino solo "1."
\makeatletter
\@ifclassloaded{book}{
  \renewcommand{\chaptermark}[1]{\markboth{\thechapter.\ #1}{}} % En el encabezado
    
  \renewcommand{\@makechapterhead}[1]{%
  \vspace*{50\p@}%
  {\parindent \z@ \raggedright \normalfont
    \ifnum \c@secnumdepth >\m@ne
      \huge\bfseries \thechapter.\hspace{1em}\ignorespaces
    \fi
    \interlinepenalty\@M
    \Huge \bfseries #1\par\nobreak
    \vskip 40\p@
  }}
}
\makeatother

%------------------------------------------------------------------------
% Paquetes de cógido
\usepackage{minted}
\renewcommand\listingscaption{Código fuente}

\usepackage{fancyvrb}
% Personaliza el tamaño de los números de línea
\renewcommand{\theFancyVerbLine}{\small\arabic{FancyVerbLine}}

% Estilo para C++
\newminted{cpp}{
    frame=lines,
    framesep=2mm,
    baselinestretch=1.2,
    linenos,
    escapeinside=||
}

% para minted
\definecolor{LightGray}{rgb}{0.95,0.95,0.92}
\setminted{
    linenos=true,
    stepnumber=5,
    numberfirstline=true,
    autogobble,
    breaklines=true,
    breakautoindent=true,
    breaksymbolleft=,
    breaksymbolright=,
    breaksymbolindentleft=0pt,
    breaksymbolindentright=0pt,
    breaksymbolsepleft=0pt,
    breaksymbolsepright=0pt,
    fontsize=\footnotesize,
    bgcolor=LightGray,
    numbersep=10pt
}


\usepackage{listings} % Para incluir código desde un archivo

\renewcommand\lstlistingname{Código Fuente}
\renewcommand\lstlistlistingname{Índice de Códigos Fuente}

% Definir colores
\definecolor{vscodepurple}{rgb}{0.5,0,0.5}
\definecolor{vscodeblue}{rgb}{0,0,0.8}
\definecolor{vscodegreen}{rgb}{0,0.5,0}
\definecolor{vscodegray}{rgb}{0.5,0.5,0.5}
\definecolor{vscodebackground}{rgb}{0.97,0.97,0.97}
\definecolor{vscodelightgray}{rgb}{0.9,0.9,0.9}

% Configuración para el estilo de C similar a VSCode
\lstdefinestyle{vscode_C}{
  backgroundcolor=\color{vscodebackground},
  commentstyle=\color{vscodegreen},
  keywordstyle=\color{vscodeblue},
  numberstyle=\tiny\color{vscodegray},
  stringstyle=\color{vscodepurple},
  basicstyle=\scriptsize\ttfamily,
  breakatwhitespace=false,
  breaklines=true,
  captionpos=b,
  keepspaces=true,
  numbers=left,
  numbersep=5pt,
  showspaces=false,
  showstringspaces=false,
  showtabs=false,
  tabsize=2,
  frame=tb,
  framerule=0pt,
  aboveskip=10pt,
  belowskip=10pt,
  xleftmargin=10pt,
  xrightmargin=10pt,
  framexleftmargin=10pt,
  framexrightmargin=10pt,
  framesep=0pt,
  rulecolor=\color{vscodelightgray},
  backgroundcolor=\color{vscodebackground},
}

%------------------------------------------------------------------------
% La etiqueda IA sirve para indicar a partir de donde le pasamos el preambulo a la IA para el escaneo
% IA

% Comandos definidos
\newcommand{\bb}[1]{\mathbb{#1}}
\newcommand{\cc}[1]{\mathcal{#1}}

% I prefer the slanted \leq
\let\oldleq\leq % save them in case they're every wanted
\let\oldgeq\geq
\renewcommand{\leq}{\leqslant}
\renewcommand{\geq}{\geqslant}

% Si y solo si
\newcommand{\sii}{\iff}

% MCD y MCM
\DeclareMathOperator{\mcd}{mcd}
\DeclareMathOperator{\mcm}{mcm}

% Signo
\DeclareMathOperator{\sgn}{sgn}

% Imagen
\DeclareMathOperator{\Ima}{Im}

% Aplicación identidad
\newcommand{\id}{\mathrm{Id}}

% Letras griegas
\newcommand{\eps}{\epsilon}
\newcommand{\veps}{\varepsilon}
\newcommand{\lm}{\lambda}

\newcommand{\ol}{\overline}
\newcommand{\ul}{\underline}
\newcommand{\wt}{\widetilde}
\newcommand{\wh}{\widehat}

\let\oldvec\vec
\renewcommand{\vec}{\overrightarrow}

% Derivadas parciales
\newcommand{\del}[2]{\frac{\partial #1}{\partial #2}}
\newcommand{\Del}[3]{\frac{\partial^{#1} #2}{\partial #3^{#1}}}
\newcommand{\deld}[2]{\dfrac{\partial #1}{\partial #2}}
\newcommand{\Deld}[3]{\dfrac{\partial^{#1} #2}{\partial #3^{#1}}}


\newcommand{\AstIg}{\stackrel{(\ast)}{=}}
\newcommand{\Hop}{\stackrel{L'H\hat{o}pital}{=}}

\newcommand{\red}[1]{{\color{red}#1}} % Para integrales, destacar los cambios.

% Método de integración
\newcommand{\MetInt}[2]{
    \left[\begin{array}{c}
        #1 \\ #2
    \end{array}\right]
}

% Declarar aplicaciones
% 1. Nombre aplicación
% 2. Dominio
% 3. Codominio
% 4. Variable
% 5. Imagen de la variable
\newcommand{\Func}[5]{
    \begin{equation*}
        \begin{array}{rrll}
            \displaystyle #1:& \displaystyle  #2 & \longrightarrow & \displaystyle  #3\\
               & \displaystyle  #4 & \longmapsto & \displaystyle  #5
        \end{array}
    \end{equation*}
}

%------------------------------------------------------------------------


\newcommand{\myhref}[2]{\href{#1}{\color{blue}\uline{#2}}}

\DeclareMathOperator{\Irr}{Irr}
\DeclareMathOperator{\Subgr}{Subgr}
\DeclareMathOperator{\Subex}{Subex}
\DeclareMathOperator{\car}{car}
\DeclareMathOperator{\Aut}{Aut}
\DeclareMathOperator{\Disc}{Disc}
\DeclareMathOperator{\Sim}{Sim}

\begin{document}

    % 1. Foto de fondo
    % 2. Título
    % 3. Encabezado Izquierdo
    % 4. Color de fondo
    % 5. Coord x del titulo
    % 6. Coord y del titulo
    % 7. Fecha

    
    % 1. Foto de fondo
% 2. Título
% 3. Encabezado Izquierdo
% 4. Color de fondo
% 5. Coord x del titulo
% 6. Coord y del titulo
% 7. Fecha
% 8. Autor

\newcommand{\portada}[8]{

    \portadaBase{#1}{#2}{#3}{#4}{#5}{#6}{#7}{#8}
    \portadaBook{#1}{#2}{#3}{#4}{#5}{#6}{#7}{#8}
}

\newcommand{\portadaExamen}[8]{

    \portadaBase{#1}{#2}{#3}{#4}{#5}{#6}{#7}{#8}
    \portadaArticle{#1}{#2}{#3}{#4}{#5}{#6}{#7}{#8}
}




\newcommand{\portadaBase}[8]{

    % Tiene la portada principal y la licencia Creative Commons
    
    % 1. Foto de fondo
    % 2. Título
    % 3. Encabezado Izquierdo
    % 4. Color de fondo
    % 5. Coord x del titulo
    % 6. Coord y del titulo
    % 7. Fecha
    % 8. Autor
    
    
    \thispagestyle{empty}               % Sin encabezado ni pie de página
    \newgeometry{margin=0cm}        % Márgenes nulos para la primera página
    
    
    % Encabezado
    \fancyhead[L]{\helv #3}
    \fancyhead[R]{\helv \nouppercase{\leftmark}}
    
    
    \pagecolor{#4}        % Color de fondo para la portada
    
    \begin{figure}[p]
        \centering
        \transparent{0.3}           % Opacidad del 30% para la imagen
        
        \includegraphics[width=\paperwidth, keepaspectratio]{../../_assets/#1}
    
        \begin{tikzpicture}[remember picture, overlay]
            \node[anchor=north west, text=white, opacity=1, font=\fontsize{60}{90}\selectfont\bfseries\sffamily, align=left] at (#5, #6) {#2};
            
            \node[anchor=south east, text=white, opacity=1, font=\fontsize{12}{18}\selectfont\sffamily, align=right] at (9.7, 3) {\href{https://losdeldgiim.github.io/}{\textbf{Los Del DGIIM}, \texttt{\footnotesize losdeldgiim.github.io}}};
            
            \node[anchor=south east, text=white, opacity=1, font=\fontsize{12}{15}\selectfont\sffamily, align=right] at (9.7, 1.8) {Doble Grado en Ingeniería Informática y Matemáticas\\Universidad de Granada};
        \end{tikzpicture}
    \end{figure}
    
    
    \restoregeometry        % Restaurar márgenes normales para las páginas subsiguientes
    \pagecolor{white}       % Restaurar el color de página
    
    
    \newpage
    \thispagestyle{empty}               % Sin encabezado ni pie de página
    \begin{tikzpicture}[remember picture, overlay]
        \node[anchor=south west, inner sep=3cm] at (current page.south west) {
            \begin{minipage}{0.5\paperwidth}
                \href{https://creativecommons.org/licenses/by-nc-nd/4.0/}{
                    \includegraphics[height=2cm]{../../_assets/Licencia.png}
                }\vspace{1cm}\\
                Esta obra está bajo una
                \href{https://creativecommons.org/licenses/by-nc-nd/4.0/}{
                    Licencia Creative Commons Atribución-NoComercial-SinDerivadas 4.0 Internacional (CC BY-NC-ND 4.0).
                }\\
    
                Eres libre de compartir y redistribuir el contenido de esta obra en cualquier medio o formato, siempre y cuando des el crédito adecuado a los autores originales y no persigas fines comerciales. 
            \end{minipage}
        };
    \end{tikzpicture}
    
    
    
    % 1. Foto de fondo
    % 2. Título
    % 3. Encabezado Izquierdo
    % 4. Color de fondo
    % 5. Coord x del titulo
    % 6. Coord y del titulo
    % 7. Fecha
    % 8. Autor


}


\newcommand{\portadaBook}[8]{

    % 1. Foto de fondo
    % 2. Título
    % 3. Encabezado Izquierdo
    % 4. Color de fondo
    % 5. Coord x del titulo
    % 6. Coord y del titulo
    % 7. Fecha
    % 8. Autor

    % Personaliza el formato del título
    \pretitle{\begin{center}\bfseries\fontsize{42}{56}\selectfont}
    \posttitle{\par\end{center}\vspace{2em}}
    
    % Personaliza el formato del autor
    \preauthor{\begin{center}\Large}
    \postauthor{\par\end{center}\vfill}
    
    % Personaliza el formato de la fecha
    \predate{\begin{center}\huge}
    \postdate{\par\end{center}\vspace{2em}}
    
    \title{#2}
    \author{\href{https://losdeldgiim.github.io/}{Los Del DGIIM, \texttt{\large losdeldgiim.github.io}}
    \\ \vspace{0.5cm}#8}
    \date{Granada, #7}
    \maketitle
    
    \tableofcontents
}




\newcommand{\portadaArticle}[8]{

    % 1. Foto de fondo
    % 2. Título
    % 3. Encabezado Izquierdo
    % 4. Color de fondo
    % 5. Coord x del titulo
    % 6. Coord y del titulo
    % 7. Fecha
    % 8. Autor

    % Personaliza el formato del título
    \pretitle{\begin{center}\bfseries\fontsize{42}{56}\selectfont}
    \posttitle{\par\end{center}\vspace{2em}}
    
    % Personaliza el formato del autor
    \preauthor{\begin{center}\Large}
    \postauthor{\par\end{center}\vspace{3em}}
    
    % Personaliza el formato de la fecha
    \predate{\begin{center}\huge}
    \postdate{\par\end{center}\vspace{5em}}
    
    \title{#2}
    \author{\href{https://losdeldgiim.github.io/}{Los Del DGIIM, \texttt{\large losdeldgiim.github.io}}
    \\ \vspace{0.5cm}#8}
    \date{Granada, #7}
    \thispagestyle{empty}               % Sin encabezado ni pie de página
    \maketitle
    \vfill
}
    \portadaExamen{ffccA4.jpg}{Álgebra III\\Examen X}{Álgebra III. Examen X}{MidnightBlue}{-8}{28}{2026}{Arturo Olivares Martos}

    \begin{description}
        \item[Asignatura] Álgebra III.
        \item[Curso Académico] 2026/27.
        \item[Grado] Doble Grado en Ingeniería Informática y Matemáticas.
        \item[Grupo] Único.
        \item[Profesor] José Gómez Torrecillas.
        \item[Descripción] Primer examen sorpresa.
        \item[Fecha] 16 de Octubre de 2025.
        \item[Duración] Una hora.
    
    \end{description}
    \newpage


    % ------------------------------------


\begin{ejercicio}
    Sea $\alpha=\sqrt[3]{2}+\sqrt{3}$. Calcular:
    \begin{enumerate}
        \item $\Irr(\alpha,\mathbb{Q})$.
        \item $\Irr(\alpha,\mathbb{Q}(\sqrt{3}))$.
        \item $\Irr(\sqrt[3]{2},\mathbb{Q}(\sqrt{3}))$.
    \end{enumerate}
\end{ejercicio}


\setcounter{ejercicio}{0}
\newpage


\begin{ejercicio}
    Sea $\alpha=\sqrt[3]{2}+\sqrt{3}$. Calcular:
    \begin{enumerate}
        \item $\Irr(\alpha,\mathbb{Q})$.
        
        Veamos en primer lugar que $\bb{Q}(\alpha) = \bb{Q}(\sqrt[3]{2},\sqrt{3})$. Por un lado:
        \begin{align*}
            \alpha-\sqrt{3} = \sqrt[3]{2} &\Longrightarrow
            \alpha^3 - 3\alpha^2\sqrt{3} + 9\alpha - 3\sqrt{3} = 2\\
            &\Longrightarrow \dfrac{\alpha^3 + 9\alpha - 2}{3(\alpha+1)} = \sqrt{3} \in \bb{Q}(\alpha)
        \end{align*}

        Por otro lado, $\sqrt[3]{2} = \alpha - \sqrt{3} \in \bb{Q}(\alpha)$, luego $\bb{Q}(\alpha) = \bb{Q}(\sqrt[3]{2},\sqrt{3})$. Buscamos un polinomio mónico en $\bb{Q}[x]$ del que $\alpha$ sea raíz. Para ello, elevamos al cuadrado la expresión de $\sqrt{3}$:
        \begin{align*}
            \dfrac{\alpha^6+2\alpha^3(9\alpha-2)+81\alpha^2+4-36\alpha}{9(\alpha^2+2\alpha+1)} &= 3 \\
            \alpha^6 + 18\alpha^4 - 4\alpha^3 + 81\alpha^2 - 36\alpha + 4 &= 27(\alpha^2+2\alpha+1) \\
            \alpha^6 + 18\alpha^4 - 4\alpha^3 + 54\alpha^2 - 90\alpha - 23 &= 0
        \end{align*}

        Por tanto, $\alpha$ es raíz de $f = x^6 + 18x^4 - 4x^3 + 54x^2 - 90x - 23\in \bb{Q}[x]$. No obstante, no es fácil ver si es irreducible o no. Para ello, consideramos las siguientes torres:
        \begin{itemize}
            \item $\bb{Q} \leq \bb{Q}(\sqrt{3}) \leq \bb{Q}(\sqrt[3]{2},\sqrt{3}) = \bb{Q}(\alpha)$
            
            Por el Lema de la Torre, tenemos que:
            \begin{equation*}
                [\bb{Q}(\alpha):\bb{Q}] = [\bb{Q}(\alpha):\bb{Q}(\sqrt{3})][\bb{Q}(\sqrt{3}):\bb{Q}]
            \end{equation*}
            \begin{itemize}
                \item $\Irr(\sqrt{3},\mathbb{Q}) = x^2-3$, que es irreducible por Eisenstein para $p=3$, luego $[\bb{Q}(\sqrt{3}):\bb{Q}] = 2$.
                \item Dado $g=x^3-2\in \bb{Q}(\sqrt{3})[x]$, tenemos que $g(\sqrt[3]{2}) = 0$, luego $\Irr(\sqrt[3]{2},\bb{Q}(\sqrt{3}))$ divide a $g$, con lo que $[\bb{Q}(\alpha):\bb{Q}(\sqrt{3})] \leq 3$. Notemos que no es fácil ver si es irreducible o no, por lo que a priori no podemos afirmar si es $1$, $2$ o $3$.
            \end{itemize}
            \item $\bb{Q} \leq \bb{Q}(\sqrt[3]{2}) \leq \bb{Q}(\sqrt[3]{2},\sqrt{3}) = \bb{Q}(\alpha)$
            
            Por el Lema de la Torre, tenemos que:
            \begin{equation*}
                [\bb{Q}(\alpha):\bb{Q}] = [\bb{Q}(\alpha):\bb{Q}(\sqrt[3]{2})][\bb{Q}(\sqrt[3]{2}):\bb{Q}]
            \end{equation*}
            \begin{itemize}
                \item $\Irr(\sqrt[3]{2},\mathbb{Q}) = x^3-2$, que es irreducible por Eisenstein para $p=2$, luego $[\bb{Q}(\sqrt[3]{2}):\bb{Q}] = 3$.
                \item Dado $h=x^2-3\in \bb{Q}(\sqrt[3]{2})[x]$, tenemos que $h(\sqrt{3}) = 0$, luego $\Irr(\sqrt{3},\bb{Q}(\sqrt[3]{2}))$ divide a $h$, con lo que $[\bb{Q}(\alpha):\bb{Q}(\sqrt[3]{2})] \leq 2$. Notemos que no es fácil ver si es irreducible o no, por lo que a priori no podemos afirmar si es $1$ o $2$.
            \end{itemize}
        \end{itemize}

        Por tanto, tenemos que $[\bb{Q}(\alpha):\bb{Q}]$ es múltiplo de $2$ y de $3$, y menor o igual que $6$. Por tanto, $[\bb{Q}(\alpha):\bb{Q}] = 6$, y $\Irr(\alpha,\mathbb{Q})$ es de grado 6. Por tanto:
        \begin{equation*}
            \Irr(\alpha,\mathbb{Q}) = x^6 + 18x^4 - 4x^3 + 54x^2 - 90x - 23\in \bb{Q}[x]
        \end{equation*}
        \item $\Irr(\alpha,\mathbb{Q}(\sqrt{3}))$.
        
        Buscamos en $\bb{Q}(\sqrt{3})[x]$ un polinomio mónico del que $\alpha$ sea raíz. De la expresión de $\sqrt{3}$ en función de $\alpha$ obtenemos que:
        \begin{equation*}
            \alpha^3+9\alpha-2 = 3(\alpha+1)\sqrt{3} \Longrightarrow \alpha^3+(9-3\sqrt{3})\alpha-2-3\sqrt{3} = 0
        \end{equation*}

        No obstante, no es fácil ver si es irreducible o no. Por la primera torre descrita, tenemos que:
        \begin{equation*}
            [\bb{Q}(\alpha):\bb{Q}(\sqrt{3})] = \dfrac{[\bb{Q}(\alpha):\bb{Q}]}{[\bb{Q}(\sqrt{3}):\bb{Q}]} = \dfrac{6}{2} = 3
        \end{equation*}

        Por tanto, tenemos que el grado de $\Irr(\alpha,\mathbb{Q}(\sqrt{3})(\alpha))$ es 3, y como se tiene que $\bb{Q}(\sqrt{3})=\bb{Q}(\sqrt{3})(\alpha)$, tenemos que:
        \begin{equation*}
            \deg(\Irr(\alpha,\mathbb{Q}(\sqrt{3}))) = 3
        \end{equation*}

        Por tanto, se deduce que:
        \begin{equation*}
            \Irr(\alpha,\mathbb{Q}(\sqrt{3})) = x^3+(9-3\sqrt{3})x-2-3\sqrt{3}\in \bb{Q}(\sqrt{3})[x]
        \end{equation*}


        \item $\Irr(\sqrt[3]{2},\mathbb{Q}(\sqrt{3}))$.
        
        Buscamos en $\bb{Q}(\sqrt{3})[x]$ un polinomio mónico del que $\sqrt[3]{2}$ sea raíz. Consideramos:
        \begin{equation*}
            x^3-2\in \bb{Q}(\sqrt{3})[x]
        \end{equation*}

        No obstante, no es fácil ver si es irreducible o no.        
        Por la primera torre descrita, tenemos que:
        \begin{equation*}
            [\bb{Q}(\alpha):\bb{Q}(\sqrt{3})] = \dfrac{[\bb{Q}(\alpha):\bb{Q}]}{[\bb{Q}(\sqrt{3}):\bb{Q}]} = \dfrac{6}{2} = 3
        \end{equation*}

        Por tanto, tenemos que el grado de $\Irr(\sqrt[3]{2},\mathbb{Q}(\sqrt{3}))$ es $[\bb{Q}(\sqrt[3]{2},\sqrt{3}):\bb{Q}(\sqrt{3})] = 3$, y como se tiene que $\bb{Q}(\sqrt[3]{2},\sqrt{3})=\bb{Q}(\alpha)$, tenemos que:
        \begin{equation*}
            \deg(\Irr(\sqrt[3]{2},\mathbb{Q}(\sqrt{3}))) = 3
        \end{equation*}

        Por tanto, se deduce que:
        \begin{equation*}
            \Irr(\sqrt[3]{2},\mathbb{Q}(\sqrt{3})) = x^3-2\in \bb{Q}(\sqrt{3})[x]
        \end{equation*}
    \end{enumerate}
\end{ejercicio}


\end{document}
